
\documentclass[12pt, b5paper]{article}
\setlength{\topmargin}{-.5in}
\setlength{\textwidth}{6.5in}
\setlength{\textheight}{9.5in}
\setlength{\oddsidemargin}{-.25in}

\usepackage{amsmath}

\newcommand{\floor}[1]{\left\lfloor #1 \right\rfloor}

\pagestyle{empty}
\begin{document}

\begin{center}
{\bf CSC 226
Exercise 1}\\
An Nguyen, V01073373\\
17 Jan 2026
\end{center}

\section*{Problem:}
Consider the recurrence T(n) below, which is defined for all integers $n \geq 0$.

$$
T(n) = 
\begin{cases}
    1 & \text{if } n = 0 \text{ or } n = 1 \\
    T(n-1) + 2T(n-2) & \text{if } n \geq 2
\end{cases}
$$
Prove that $T(n) \leq 2^n$ for all $n \geq 0$.

\section*{Proof:}

\textbf{Idea: }
We solve the recurrence relation using proof by \textbf{strong induction} on $n$. \\
\newline
\textbf{Base Cases: }
\begin{align*}
\text{For } n = 0: \quad T(0) &= 1 = 2^0 = 1 \\
\text{For } n = 1: \quad T(1) &= 1 < 2^1 = 2
\end{align*}
Therefore, our base cases hold. \\
\newline
\textbf{Induction Hypothesis: }
Suppose that $T(j) \leq 2^j$ for all $0 \leq j \leq k$, where $k \geq 1$.\\
\newline
\textbf{Induction Step: } \\
We must show that the statement holds for $n = k + 1$: 
$$T(k+1) \leq 2^{k+1}$$
Using the recurrence relation, we have:
\begin{align*}
T(k+1) &= T(k) + 2T(k-1) \\
&\leq 2^k + 2 \cdot 2^{k-1}
= 2^k + 2^k = 2 \cdot 2^k = 2^{k+1}  \text{ (by I.H.)} 
\end{align*}
\textbf{Conclusion: } \\
Therefore, by induction, the claim holds for all $n \geq 0$.

\newpage
\begin{center}
{\bf CSC 226
Exercise 2}\\
An Nguyen, V01073373\\
24 Jan 2026
\end{center}

\section*{Problem:}
We have seen that the Quickselect algorithm has worst-case $\theta(n)$ performance if pivots are 
selected with the BFPRT pivot selection scheme using sub-arrays of size $5$. Prove that Quickselect 
has worst-case $\theta(n)$ performance if pivots are selected with the BFPRT pivot selection scheme 
using sub-arrays of size $7$.
\section*{Proof:}
\subsection*{Idea: }
We'll compute the largest size of the two partitions created by using BFPRT algorithm with sub-arrays of size $7$. 
Then, we'll set up the recurrence relation for Quickselect and solve it using induction on $n$. 
\subsection*{Calculating the size of the partitions: }
For an array of size n, we have $\floor{\frac{n}{7}}$ sub-arrays of size $7$ (ignore the 
remaining elements). After the partitioning is complete, each sub-array is sorted in $\theta(1)$ time 
since the size is constant. Since the sub-array size is odd, each sub-array will have a median 
after sorting.\\
\newline
We find the median of each sub-array, which gives us $\frac{n}{7}$ medians. The time taken to sort 
all sub-arrays and find their medians is $\theta(n)$. Suppose we arrange the sub-arrays in 
non-decreasing order based on their medians. Let $p$ be the median of the middle sub-array. This value 
is the median of the sequence of medians, so we call it \textbf{median of medians}. \\
\newline
Since there are $\frac{n}{7}$ sub-arrays, $p$ is greater than the $\floor{\frac{n}{14}}$ medians and 
less than the $\floor{\frac{n}{14}}$ medians. Each median is greater than $3$ elements in its sub-array 
and less than $3$ elements in its sub-array, so:
$$|L| \geq \frac{n}{14} \times 4 = \frac{4n}{14}$$
$$|G| \geq \frac{n}{14} \times 4 = \frac{4n}{14}$$\\
($L$ is the set of elements $\leq p$, and $G$ is the set of elements $\geq p$.) \\
\newline
Since $|L| + |G| + 1 = n$, the size of the largest partition is at most:
$$n - \frac{4n}{14} = \frac{10n}{14} = \frac{5n}{7}$$
Therefore, recursing on the largest partition takes at most $T\left(\frac{5n}{7}\right)$ time. \\
\newline
Computing the sequence of sub-array medians take $\theta(n)$ time, and finding $p$ takes a recursive 
call to Quickselect on the sequence of
$\frac{n}{7}$ medians, which takes $T\left(\frac{n}{7}\right)$ time.
\subsection*{Recurrence:}
The overall running time of Quickselect with BFPRT pivoting can be modelled by the recurrence:
$$
T(n) = 
\begin{cases}
    c_b & \text{if } n \leq 1 \\
    T\left(\floor{\frac{5n}{7}}\right) + T\left(\floor{\frac{n}{7}}\right) + c_1n + c_2 & \text{if } n \geq 2
\end{cases}
$$
Notice that $\frac{5n}{7} + \frac{n}{7} = \frac{6n}{7} < n$ , so the size of the subproblems decreases 
geometrically. \\
Let $c=c_b + c_1 + c_2$ and note that since $c_{1}n + c_{2}n$ for $n\geq 2$, the simplified recurrence below is 
always greate than $T(n)$:
$$
T'(n) = 
\begin{cases}
    c & \text{if } n \leq 1 \\
    T'\left(\lfloor \frac{5n}{7} \rfloor\right) + T'\left(\lfloor\frac{n}{7}\rfloor\right) + cn & \text{if } n \geq 2
\end{cases}
$$
\newline
\textbf{Claim: } $T'(n) \leq 7cn$ for $n \geq 1$. \\
\newline
\textbf{Proof: } By induction on $n$. \\
\newline
\textbf{Base Case: }
When $n=1$, $T'(1) = c = cn \leq 7cn$. So the claim holds. \\
\newline
\textbf{Induction Hypothesis: }
Suppose $T'(n) \leq 7cn$ for $n \in {1, 2, ..., k}$ for some $k \geq 1$. \\
\newline
\textbf{Induction Step: }
Consider $T'(n + 1)$:
\begin{align*}  
T'(n+1) &= T'\left(\floor{\frac{5(n+1)}{7}}\right) + T'\left(\floor{\frac{n+1}{7}}\right) + c(n+1) \\
&\leq 7c \cdot \left(\floor{\frac{5(n+1)}{7}}\right) + 7c \cdot \left(\floor{\frac{n+1}{7}}\right) + c(n+1) \text{ (by I.H.)} \\
&\leq 7c \cdot \left(\frac{5(n+1)}{7}\right) + 7c \cdot \left(\frac{n+1}{7}\right) + c(n+1) \\
&\leq 5c(n+1) + c(n+1) + c(n+1) \\
&= 7cn
\end{align*}
Therefore, the claim holds for $n+1$ and, by induction, the claim holds for all $n\geq1$.

\end{document}